% This must be in the first 6 lines to tell arXiv to use pdfLaTeX, which is strongly recommended.
\pdfoutput=1
% In particular, the hyperref package requires pdfLaTeX in order to break URLs across lines.

\documentclass[11pt]{article}

% Remove the "review" option to generate the final version.
\usepackage{acl}

% Standard package includes
\usepackage{times}
\usepackage{latexsym}
\usepackage{url}

% For proper rendering and hyphenation of words containing Latin characters (including in bib files)
\usepackage[T1]{fontenc}
% For Vietnamese characters
% \usepackage[T5]{fontenc}
% See https://www.latex-project.org/help/documentation/encguide.pdf for other character sets

% This assumes your files are encoded as UTF8
\usepackage[utf8]{inputenc}

% This is not strictly necessary, and may be commented out,
% but it will improve the layout of the manuscript,
% and will typically save some space.
\usepackage{microtype}

% If the title and author information does not fit in the area allocated, uncomment the following
%
\setlength\titlebox{200pt}
%
% and set <dim> to something 5cm or larger.

\title{COLING 2025 (Non-Mandatory) Workshop Proposal Template}

% Author information can be set in various styles:
% For several authors from the same institution:
% \author{Author 1 \and ... \and Author n \\
%         Address line \\ ... \\ Address line}
% if the names do not fit well on one line use
%         Author 1 \\ {\bf Author 2} \\ ... \\ {\bf Author n} \\
% For authors from different institutions:
% \author{Author 1 \\ Address line \\  ... \\ Address line
%         \And  ... \And
%         Author n \\ Address line \\ ... \\ Address line}
% To start a separate ``row'' of authors use \AND, as in
% \author{Author 1 \\ Address line \\  ... \\ Address line
%         \AND
%         Author 2 \\ Address line \\ ... \\ Address line \And
%         Author 3 \\ Address line \\ ... \\ Address line}

\author{ Mo El-Haj \and Katsuhito Sudoh \\
COLING 2025 Workshop Co-Chairs \\
\texttt{foobar@coling2025.org}
}

\begin{document}
\maketitle
\begin{abstract}
The International Conference on Computational Linguistics (COLING) invites proposals for workshops to be held in conjunction with the 31st edition of COLING in 2025  in Abu Dhabi (UAE) \cite{coling2025}.
We solicit proposals in all areas of computational linguistics, language resources and evaluation, broadly conceived to include related disciplines such as linguistics, language documentation, natural language processing, speech and multimodal processing, computational social science, and the digital humanities.

Please note that for this edition of COLING, \textbf{we will not be conducting a joint call for workshops with the ACL conferences (AACL, ACL, EACL, EMNLP, NAACL)}. This decision is due to scheduling differences between COLING and the ACL conferences in 2025. COLING is happening much earlier than the *ACL conferences, which makes it impractical to coordinate a joint call. However, we anticipate  a return to joint workshop calls in the future.
\end{abstract}

\section{SUBMISSION INFORMATION}
Proposals should be submitted as PDF documents. Note that submissions should essentially be ready to be turned into a Call for Workshop Papers shortly after the  notification.

The proposals should be at most two pages long for the main proposal, with additional two pages permitted for information about organisers, tentative program committee, and references. Hence, the entire proposal should not exceed FOUR pages in length, excluding references. Please note that workshops should either be 100\% in-person or 100\% virtual; hybrid formats will not be allowed. For in-person workshops, at least one workshop organiser should be physically present to run the workshop at COLING.


\paragraph{The two pages for the main proposal must include:}
\begin{itemize}
    \item A title and a brief description of the workshop topic and content.
    \item A list of invited speakers, if applicable, with an indication of which ones have already agreed and which are tentative, and sources of funding for the speakers, if needed.
    \item An estimate of the number of attendees.
    \item Workshop format: in-person or virtual (hybrid proposals will not be accepted).
    \item A description of any shared tasks associated with the workshop, and estimate of the number of participants. Note that any shared task will also need to be reviewed by the workshop committee for ethical concerns.
    \item A description of special requirements and technical needs, where relevant.
    \item If the workshop has been held before, a note specifying where previous iterations of the workshops were held, how many submissions the workshop received, how many papers were accepted (also specify whether they were not regular papers, e.g., shared task system description papers, non-archival papers), and how many attendees the workshop attracted.

\end{itemize}

\paragraph{The two pages for information about organisers, program committee, and references must include:}
\begin{itemize}
    \item The names, affiliations, and email addresses of the organisers, with one-paragraph statements of their research interests, areas of expertise, and experience in organising workshops and related events.
    \item A list of Program Committee members, with an indication of which members have already agreed. Organisers should do their best to estimate the number of submissions (especially for recurring workshops) in order to (a) ensure a sufficient number of reviewers so that each paper receives 3 reviews, and (b) anticipate that no one is committed to reviewing more than 3 papers. This practice is likely to ensure on-time, and more thorough and thoughtful reviews.
    \item A very brief advertisement or tagline for the workshop, up to 140 characters, that highlights any key information you wish prospective attendees to know, and which would be suitable to be put onto a web-based survey (see below).
    \item Unlimited space for references

\end{itemize}

\paragraph{The proposals should be submitted no later than June 14, 2024, 11:59 PM Samoa Standard Time (SST) (UTC/GMT-11, ‘Anywhere on Earth’).}

\paragraph{Submission is electronic, and the submission URL will be displayed on the COLING 2025 website: \url{https://coling2025.org/}.}

\paragraph{The workshop proposals will be evaluated according to their originality and impact, the expected interest level of participants, as well as the quality of the organising team and Program Committee and their contribution to the diversity of the conference.}

\section{IMPORTANT DATES}
\begin{itemize}
    \item 14 June 2024: Proposal submission due
    \item 30 June 2024: Notification of acceptance
    \item 19-20 January 2025: COLING-2025 Workshops
    \item 21-24 January 2025: COLING-2025 conference
\end{itemize}

\section{DIVERSITY AND INCLUSION}
We particularly encourage submissions of underrepresented groups in computational linguistics, language resource and evaluation, including researchers from any demographic or geographic minority, with disabilities, or others. In the evaluation of the proposal, we will take these aspects into account to create a varied and balanced set of tutorials.

This includes several aspects of diversity, namely (1) how the topic of the tutorial contributes to improved diversity and increased fairness in the field, (2) if the topic is particularly relevant for a specific underrepresented group of potential participants, (3), if the presenters are from an underrepresented group.


\section{WORKSHOP ORGANIZER RESPONSIBILITIES}
The organisers of the accepted workshops must be present in person in Abu Dhabi for an in-person workshop and virtually for a virtual workshop. They will be responsible for publicising and conducting the workshop, which includes reviewing submissions, preparing the camera-ready workshop proceedings, coordinating the meeting days, and ensuring that all participants are informed about COLING’s anti-harassment policy. (see \url{https://coling2022.org/policy}). It is crucial that organisers commit to all deadlines. In particular, failure to produce the camera-ready proceedings on time will lead to the exclusion of the workshop from the unified proceedings and author indexes. Workshop organisers cannot accept submissions for publication that will be (or have been) published elsewhere, although they are free to set their own policies on simultaneous submission and review, as well as to accept additional non-archival presentations.


% Entries for the entire Anthology, followed by custom entries
\bibliography{custom}

\end{document}
